\documentclass[10pt,a4paper]{amsart}
\usepackage{amsmath,amssymb,amsthm,graphicx}
% Retained preamble for the Stacks Project chapter snapshot in this directory.
% The chapter's own preamble.tex is a copy of the shared Stacks Project preamble; the block
% below is sliced MECHANICALLY, line for line, from it. Nothing was added, removed or
% redefined; the wrapper class, the xr-hyper/hyperref wiring and the externaldocument list
% are omitted because the assembler supplies the document class.
%
% lines taken from preamble.tex: 3, 5, 9, 11, 12, 156, 157, 158, ... (51 lines total)
%
% package coverage restored from the chapter's own preamble (2026-10-13): xy, xy-2cell

\usepackage[all]{xy}
\xyoption{2cell}
\UseAllTwocells

\newenvironment{reference}{\comment}{\endcomment}
\newenvironment{slogan}{\comment}{\endcomment}
\newenvironment{history}{\comment}{\endcomment}
\theoremstyle{plain}
\newtheorem{theorem}[subsection]{Theorem}
\newtheorem{proposition}[subsection]{Proposition}
\newtheorem{lemma}[subsection]{Lemma}
\theoremstyle{definition}
\newtheorem{definition}[subsection]{Definition}
\newtheorem{example}[subsection]{Example}
\newtheorem{exercise}[subsection]{Exercise}
\newtheorem{situation}[subsection]{Situation}
\theoremstyle{remark}
\newtheorem{remark}[subsection]{Remark}
\newtheorem{remarks}[subsection]{Remarks}
\numberwithin{equation}{subsection}
\def\lim{\mathop{\mathrm{lim}}\nolimits}
\def\colim{\mathop{\mathrm{colim}}\nolimits}
\def\Spec{\mathop{\mathrm{Spec}}}
\def\Hom{\mathop{\mathrm{Hom}}\nolimits}
\def\Ext{\mathop{\mathrm{Ext}}\nolimits}
\def\SheafHom{\mathop{\mathcal{H}\!\mathit{om}}\nolimits}
\def\SheafExt{\mathop{\mathcal{E}\!\mathit{xt}}\nolimits}
\def\Sch{\mathit{Sch}}
\def\Mor{\mathop{\mathrm{Mor}}\nolimits}
\def\Ob{\mathop{\mathrm{Ob}}\nolimits}
\def\Sh{\mathop{\mathit{Sh}}\nolimits}
\def\NL{\mathop{N\!L}\nolimits}
\def\CH{\mathop{\mathrm{CH}}\nolimits}
\def\proetale{{pro\text{-}\acute{e}tale}}
\def\etale{{\acute{e}tale}}
\def\QCoh{\mathit{QCoh}}
\def\Ker{\mathop{\mathrm{Ker}}}
\def\Im{\mathop{\mathrm{Im}}}
\def\Coker{\mathop{\mathrm{Coker}}}
\def\Coim{\mathop{\mathrm{Coim}}}
\DeclareMathSymbol{\boxtimes}{\mathbin}{AMSa}{"02}
\def\QCohstack{\mathcal{QC}\!\mathit{oh}}
\def\Cohstack{\mathcal{C}\!\mathit{oh}}
\def\Spacesstack{\mathcal{S}\!\mathit{paces}}
\def\Quotfunctor{\mathrm{Quot}}
\def\Hilbfunctor{\mathrm{Hilb}}
\def\Curvesstack{\mathcal{C}\!\mathit{urves}}
\def\Polarizedstack{\mathcal{P}\!\mathit{olarized}}
\def\Complexesstack{\mathcal{C}\!\mathit{omplexes}}
\def\Pic{\mathop{\mathrm{Pic}}\nolimits}
\def\Picardstack{\mathcal{P}\!\mathit{ic}}
\def\Picardfunctor{\mathrm{Pic}}
\def\Deformationcategory{\mathcal{D}\!\mathit{ef}}
\usepackage{cleveref}
\crefname{originallemma}{original lemma}{original lemmas}
\Crefname{originallemma}{Original lemma}{Original lemmas}
\makeatletter
\newlabel{lemma-resolve}{{0FSH}{source}{}{sourcepointer-1}{}}
\newlabel{lemma-resolve@cref}{{[originallemma][0][]0FSH}{[1][0][]source}{}{}{}}
\makeatother
\crefname{originallemma}{original lemma}{original lemmas}
\Crefname{originallemma}{Original lemma}{Original lemmas}
\makeatletter
\newlabel{lemma-special-good}{{0FU1}{source}{}{sourcepointer-2}{}}
\newlabel{lemma-special-good@cref}{{[originallemma][0][]0FU1}{[1][0][]source}{}{}{}}
\makeatother
\crefname{originallemma}{original lemma}{original lemmas}
\Crefname{originallemma}{Original lemma}{Original lemmas}
\makeatletter
\newlabel{remark-sheaf-graded-sets}{{0FSF}{source}{}{sourcepointer-3}{}}
\newlabel{remark-sheaf-graded-sets@cref}{{[originallemma][0][]0FSF}{[1][0][]source}{}{}{}}
\makeatother
\title{Stacks, Differential Graded Sheaves: lemma-good-dga}
\author{The Stacks Project Authors}
\date{stacks-project:sdga.tex:local-snapshot-20261012; GNU FDL 1.2 or later}
\begin{document}
\maketitle
\noindent\textit{Editorial note.} Statement span is lines 5759-5768 of sources/stacks-sdga/pinned-original-source.tex; proof span is lines 5770-5830, the complete author proof (59 body lines).\par\medskip
\noindent\textit{Editorial note.} Navigation stubs: 3 cited-but-not-included labels are supplied as non-load-bearing pointers with their pinned tags where the snapshot tag index has them.\par\medskip
\noindent\textit{Editorial note.} extra\_packages is empty on purpose: the retained preamble already loads xy with options and a second declaration clashes in the assembler wrapper.\par\medskip
\noindent\textit{Editorial note.} Text is retained verbatim, including the author's own typos and placeholders.\par\medskip
\noindent\textit{Editorial note.} Queue index 8 of queue\_sdga.json (pb=2440, stmt 5759-5768, proof 5770-5830).\par\medskip
\noindent\textit{Allowed standard prerequisites (editorial).}\begin{itemize}
\item Explicit prerequisite boundary. The cited inputs are prerequisites; the argument in the retained span is not replaced by any cited result.
\end{itemize}\medskip
\ifdefined\hypertarget\hypertarget{sourcepointer-1}{}\fi
\noindent\textit{Source reference.} Original lemma 0FSH: sources/stacks-sdga/pinned-original-source.tex, source lines 3186--3196. Cited by the retained statement or proof but not part of the collected spans; supplied only as a navigation pointer. The argument does not depend on its content.\par\medskip
\ifdefined\hypertarget\hypertarget{sourcepointer-2}{}\fi
\noindent\textit{Source reference.} Original lemma 0FU1: sources/stacks-sdga/pinned-original-source.tex, source lines 5632--5643. Cited by the retained statement or proof but not part of the collected spans; supplied only as a navigation pointer. The argument does not depend on its content.\par\medskip
\ifdefined\hypertarget\hypertarget{sourcepointer-3}{}\fi
\noindent\textit{Source reference.} Original lemma 0FSF: sources/stacks-sdga/pinned-original-source.tex, source lines 3098--3126. Cited by the retained statement or proof but not part of the collected spans; supplied only as a navigation pointer. The argument does not depend on its content.\par\medskip
\begin{lemma}
\label{lemma-good-dga}
Let $(\mathcal{C}, \mathcal{O})$ be a ringed site.
Let $(\mathcal{B}, \text{d})$ be a differential graded $\mathcal{O}$-algebra.
There exists a quasi-isomorphism of differential graded $\mathcal{O}$-algebras
$(\mathcal{A}, \text{d}) \to (\mathcal{B}, \text{d})$ such that
$\mathcal{A}$ is graded flat and K-flat as a complex of $\mathcal{O}$-modules
and such that the same is true after pullback by any morphism of
ringed topoi.
\end{lemma}

\begin{proof}
The proof is exactly the same as the first proof of
Lemma \ref{lemma-resolve} but now working with free graded
algebras instead of free graded modules.

\medskip\noindent
We will construct $\mathcal{A} = \colim \mathcal{A}_i$ as in
Lemma \ref{lemma-special-good} by constructing
$$
\mathcal{A}_0 \to \mathcal{A}_1 \to \mathcal{A}_2 \to \ldots \to \mathcal{B}
$$
Let $\mathcal{S}_0$ be the sheaf of graded sets
(Remark \ref{remark-sheaf-graded-sets})
whose degree $n$ part is $\Ker(\text{d}_\mathcal{B}^n)$.
Consider the homomorphism of differential
graded modules
$$
\mathcal{A}_0 = \mathcal{O}\langle \mathcal{S}_0 \rangle
\longrightarrow
\mathcal{B}
$$
where map sends a local section $s$ of $\mathcal{S}_0$
to the corresponding local section of $\mathcal{A}^{\deg(s)}$
(which is in the kernel of the differential, so our map is
a map of differential graded algebras indeed). By construction the
induced maps on cohomology sheaves $H^n(\mathcal{A}_0) \to H^n(\mathcal{B})$
are surjective and hence the same will remain true for all $i$.

\medskip\noindent
Induction step of the construction. Given $\mathcal{A}_i \to \mathcal{B}$
denote $\mathcal{S}_{i + 1}$ the sheaf of graded sets whose degree $n$ part is
$$
\Ker(\text{d}_{\mathcal{A}_i}^{n + 1})
\times_{\mathcal{B}^{n + 1}, \text{d}}
\mathcal{B}^n
$$
This comes equipped with a canonical map
$$
\delta_{i + 1} : \mathcal{S}_{i + 1} \longrightarrow
\mathcal{A}_i
$$
whose image is contained in the kernel of $\text{d}_{\mathcal{A}_i}$
by construction. Hence $\mathcal{A}_{i + 1} = \mathcal{O}\langle
\mathcal{S}_0 \amalg \ldots \mathcal{S}_{i + 1}\rangle$ has a differential
exteding the differential on $\mathcal{A}_i$, see discussion at the
start of this section. The map from $\mathcal{A}_{i + 1}$ to
$\mathcal{B}$ is the unique map of graded algebras which
restricts to the given map on $\mathcal{A}_i$ and sends
a local section $s = (a, b)$ of $\mathcal{S}_{i + 1}$
to $b$ in $\mathcal{B}$. This is compatible with differentials
exactly because $\text{d}(b)$ is the image of $a$ in $\mathcal{B}$.

\medskip\noindent
The map $\mathcal{A} \to \mathcal{B}$ is a quasi-isomorphism:
we have $H^n(\mathcal{A}) = \colim H^n(\mathcal{A}_i)$
and for each $i$ the map $H^n(\mathcal{A}_i) \to H^n(\mathcal{B})$
is surjective with kernel annihilated by the map
$H^n(\mathcal{A}_i) \to H^n(\mathcal{A}_{i + 1})$ by construction.
Finally, the flatness condition for $\mathcal{A}$ where shown in
Lemma \ref{lemma-special-good}.
\end{proof}

\end{document}
